\begin{examplebox}
\textbf{例1：已知切点，三秒写切线（基础母题）。}

求椭圆 $x^2/8+y^2/2=1$ 在 $P(2,1)$ 处的切线。

\textbf{解：}先核验 $4/8+1/2=1$，切点合法。直接把平方拆为“切点乘变量”：
\[
 \frac{2x}{8}+\frac{y}{2}=1
 \quad\Longrightarrow\quad
 \boxed{x+2y-4=0}.
\]
\textbf{最快识别：}题干给“椭圆上一点”，不用联立、求导、判别式。

\medskip
\textbf{例2：已知外点，求两条切线的切点连线（重要母题）。}

从 $M(2,2)$ 向椭圆 $x^2/4+y^2=1$ 引两条切线，求切点弦及两切点。

\textbf{解：} $M$ 的位置量为 $S=2^2/4+2^2=5>1$，确在椭圆外。直接写极线
\[
 \frac{2x}{4}+2y=1
 \quad\Longrightarrow\quad
 \boxed{x+4y-2=0}.
\]
为求切点，联立 $x=2-4y$ 与椭圆：
\[
 \frac{(2-4y)^2}{4}+y^2=1
 \quad\Longrightarrow\quad y(5y-4)=0.
\]
故
\[
 \boxed{T_1=(2,0),\qquad T_2=(-6/5,4/5).}
\]
\textbf{回代验算：} $T_1$ 处切线为 $x=2$，包含 $M(2,2)$；$T_2$ 处切线为 $-3x+8y=10$，代入 $M$ 得 $-6+16=10$，两条确实都是从 $M$ 出发的切线。\textbf{千万别把 $x+4y-2=0$ 误认成经过 $M$ 的切线。}

\medskip
\textbf{例3：切点在坐标轴上（专门防止斜率除零）。}

椭圆 $x^2/9+y^2/4=1$，求 $P_1=(3,0)$、$P_2=(0,-2)$ 处的切线。

\textbf{解：}切点式直接给出
\[
 P_1:\ \frac{3x}{9}=1\ \Longrightarrow\ \boxed{x=3},
 \qquad
 P_2:\ \frac{-2y}{4}=1\ \Longrightarrow\ \boxed{y=-2}.
\]
$P_1$ 处斜率不存在，但切点式完全有效；$P_2$ 处斜率为 $0$。

\medskip
\textbf{例4：反向应用，判定含参数的直线何时相切。}

求实数 $t$，使 $2x+3y=t$ 与椭圆 $x^2/9+y^2/4=1$ 相切，并求相应切点。

\textbf{解：}相切条件为
\[
 t^2=9\cdot2^2+4\cdot3^2=72
 \quad\Longrightarrow\quad
 \boxed{t=\pm6\sqrt2}.
\]
对于 $t=6\sqrt2$，切点由“系数反推切点”得到
\[
 x_0=\frac{a^2A}{t}=\frac{18}{6\sqrt2}=\frac{3}{\sqrt2},
 \qquad
 y_0=\frac{b^2B}{t}=\frac{12}{6\sqrt2}=\sqrt2.
\]
故 $P_+=(3/\sqrt2,\sqrt2)$；同理 $t=-6\sqrt2$ 时 $P_-=(-3/\sqrt2,-\sqrt2)$。两点均满足椭圆方程。\textbf{用途：}带参数的直线判相切时，比完整联立二次方程更快。

\medskip
\textbf{例5：极线并不总是“切点连线”（位置分类）。}

在椭圆 $x^2/4+y^2=1$ 上，分别研究 $M_1=(4,0)$、$M_2=(1/2,0)$、$M_3=(0,0)$。

\textbf{解：}
\begin{itemize}
    \item $M_1$ 的位置量 $S=4>1$：极线为 $4x/4=1$，即 $x=1$；与椭圆交于 $(1,\sqrt3/2)$、$(1,-\sqrt3/2)$，确为两个切点。
    \item $M_2$ 的位置量 $S=1/16<1$：形式方程为 $(1/2)x/4=x/8=1$，即 $x=8$；它与椭圆没有交点，因此\textbf{没有实切点，也谈不上实切点连线}。
    \item $M_3$ 在原点：代换式为 $0=1$，\textbf{不是一条普通实直线}。
\end{itemize}

\medskip
\textbf{高价值变式：}若把例1的切点 $P(2,1)$ 改为 $M(2,2)$，还能直接说“切线为 $2x/8+2y/2=1$”吗？

\textbf{答：不能。}因为 $2^2/8+2^2/2=5/2>1$，该点不在椭圆上；直接代换得到的是\textbf{极线}，不是 $M$ 处的切线。若题目要求从 $M$ 引出的两条切线，还须先求切点，分别写两条切点式。
\end{examplebox}
